\documentclass[10pt,a4paper,openany,twoside,prof]{book}

\usepackage{layout_cours}
\usepackage{hyperref}

\setlist{leftmargin=5mm}

\begin{document}

	\chapnumber{3} % numéro du chapitre
	\titre{Arithmétique} % titre du chapitre
	\classe{Seconde} % classe
	\entetegal

	% Debut du texte
	% **************
	\vspace{-7mm}
	\paragraph*{Objectif du chapitre:}
	\begin{itemize}
		\vspace{-2mm}
		\setlength\itemsep{-1mm}
		\item Connaître le vocabulaire de l'arithmétique: diviseur, multiple, nombre premier
		\item Savoir passer d'une notion intuitive (comme celle de nombre pair) à sa formalisation mathématique (il s'ecrit $2k$ avec $k\in\Z$)
		\item Utiliser l'arithmétique pour résoudre des problèmes
	\end{itemize}
	\vspace{-4mm}

\section{Divisibilité}

\begin{defi}{Divisibilité}{}
Un nombre entier $a\in\Z$ est \bf{divisible} par un nombre entier $b\in\Z$ s'il existe un nombre entier $k\in\Z$ tel que: $a=k\times b$.

On dit que:
\begin{itemize}
	\item $k$ est le \bf{facteur multiplicatif}
	\item $a$ est un \bf{multiple} de $b$ ou encore il est \bf{divisible} par $b$
	\item $b$ est un \bf{diviseur} de $a$
\end{itemize}
\end{defi}

\begin{ex}{}{}
$72$ est divisible par $24$ car $72 = 3\times 24$

 On peut donc dire que $72$ est un multiple de $24$ et de $3$ ou que $24$ est un diviseur de $72$.
\end{ex}

\begin{rem}
	$72$ est également divisible par $8$, $9$ et même par $-3$.
\end{rem}

\begin{methode}{Déterminer si un grand nombre $a$ est un multiple de $b$}{}
  On réalise la division euclidienne de $a$ par $b$. On trouve un quotient $q$ et un reste $r$.
  \begin{itemize}
	\item
	Si le reste est $r=0$, alors on peut écrire $a=b\times q$ et $b$ divise $a$.
	\item
	Si le reste est $r\ne 0$, alors on a $a=b\times q+r$ et $b$ ne divise pas $a$.
  \end{itemize}
\end{methode}

\begin{ex}{}{}
  On veut savoir si $a=45687$ est divisible par $b=7$. On trouve un reste non nul. Donc $a$ n'est pas divisible par $b$.
\end{ex}

\begin{defi}{Parité}{}
Un entier multiple de deux est un entier \bf{pair}, les autres sont les entiers \bf{impairs}.
\end{defi}

\begin{prop}{}{}
	Soit $a\in\Z$.

	a est pair, si et seulement s'il existe $k\in\Z$ tel que $a=2k$.

	a est impair, si et seulement s'il existe $k\in\Z$ tel que $a=2k+1$.
\end{prop}

\begin{ex}{}{}
$72$ est un multiple de $2$. En effet, $2\times 36 = 72$. Donc 72 est un entier pair.

$27$ est impair. On peut effectivement écrire $27 = 26+1 = 2\times13+1$
\end{ex}

\begin{ex}{}{}
	Soit $n=56 = 8\times7$ et $m=21=3\times7$. Ce sont deux multiples de $a= 7$.

	On observe que $n+m=77$ est aussi un multiple de $a=7$. D'ailleurs on peut expliciter le coefficient multiplicatif.
	$$n+m= 8\times7+3\times7 = (8+3)\times7=11\times7$$
\end{ex}

\begin{prop}{}{}
Soit $a\in\Z$ un entier, la somme de deux multiples de $a$ est un multiple de $a$.
\end{prop}

\begin{deme}{}{}
Soit $a\in\N$. Soient $n, m\in\N$ des multiples de $a$.

Il existe $k\in\N$ et $k'\in\N$ deux entiers tel que:
$$ n =k\times a  \quad \text{ et } \quad m = k' \times a$$

Par addition, on a:
$$ n + m = k\times a  + k' \times a $$

En factorisant par $a$, il vient:
$$ n + m = (k  + k') \times a $$

Donc $n+m$ est un multiple de $a$.
\end{deme}

\begin{ex}{}{}
	On se donne un nombre pair $a=13$. Alors son carré est :

	$a^2=13^2=(10+3)^2 = 10^2+2\times10\times3+3^2 = 169$

	Il est encore impair.
\end{ex}


\begin{prop}{}{}
Le carré d'un nombre pair est pair. Le carré d'un nombre impair est impair.
\end{prop}

\begin{deme}{}{}
	Montrons le cas d'un nombre impair $a\in\N$. Il existe $k\in\N$ tel que:
$$ a =2k+1 $$

Par multiplication, on a:
$$ a^2 = \left(2k+1 \right)^2 = (2k)^2 + 2 \times (2k) \times 1 + 1^2 = 4k^2 + 4k +1 = 2(2k^2+2k)+1$$

Donc $a^2$ est un nombre impair.

\medskip
Montrons le cas d'un nombre pair $b\in\N$. Il existe $k\in\N$ tel que:
$$ b =2k $$

On a:
$$ b^2 = \left(2k \right)^2 = 4k^2= 2\times(2k)$$

Donc $b^2$ est un nombre pair.
\end{deme}

\begin{prop}{}{}
\begin{itemize}
\item  Un nombre est divisible par $2$ si il se termine par $0$, $2$, $4$, $6$ ou $8$
\item  Un nombre est divisible par $5$ s'il se termine par $0$ ou $5$.

\item  Un nombre est divisible par $3$ si la somme de ses chiffres est divisible par $3$.


\item  Un nombre est divisible par $9$ si la somme de ses chiffres est divisible par $9$.

\end{itemize}
\end{prop}

\begin{ex}{}{}
  $2 538$ est divisible par $9$ car $2 + 5 + 3 + 8 = 18$ qui est divisible par $9$ ($3 \times 9 = 18$).

   $249$ est divisible par $3$ car $2 + 4 + 9 = 15$ qui est divisible par $3$ ($3 \times 5 = 15$).
\end{ex}


\begin{defi}{Fraction irréductible}{}
  Soit $a\in\Z$ et $b\in\N^*$. On dit que la fraction $\frac{a}{b}$ est irréductible lorsque le seul diviseur commun à $a$ et $b$ est $1$. On ne peut pas la simplifier.
\end{defi}

\begin{ex}{}{}
  $\frac{15}{6}$ n'est pas irréductible car $3$ est un diviseur commun de $15$ et de $6$.

\smallskip
  $\frac{15}{7}$ est irréductible car l'ensemble des diviseurs (positifs) de $15$ est $\{1,3,5,15\}$ et l'ensemble des diviseurs de $7$ est $\{1,7\}$.
\end{ex}

\begin{prop}{}{}
	$\sqrt{2}$ est irrationnel. Autrement dit, $\sqrt{2}\notin\Q$.
\end{prop}

\begin{deme}{}{}
	On raisonne par l'absurde, en supposant que $\sqrt{2}\in\Q$ (c'est-à-dire qu'il est rationnel). On va montrer que cela conduit à une contradiction de logique.

	Il existe $a\in\Z$ et $b\in\N^*$ des entiers tels que $\sqrt{2} = \frac{a}{b}$ et que la fraction est irréductible. Alors en élevant au carré $2 = \frac{a^2}{b^2}$. Donc $a^2=2\times b^2$. Donc $a^2$ est pair. Alors $a$ aussi est pair.

	On écrit donc $a=2k$ avec $k\in\Z$. L'égalité $a^2=2b^2$ devient $4k^2 = 2b^2$. On simplifie pour obtenir $2k^2=b^2$. Par le même raisonnement on en déduit que $b$ est pair.

	Alors la fraction $\frac{a}{b}$ se simplifie par $2$ et elle n'est pas irréductible. C'est absurde. L'hypothèse initiale est fausse et $\sqrt{2}$ n'est pas rationnel.
\end{deme}

\section{Nombres premiers}
\begin{defi}{Nombre premier}{}
Un entier naturel $n\in\N$ qui a exactement deux diviseurs distincts (à savoir $1$ et $n$) est appelé un \bf{nombre premier}.
\end{defi}

\begin{ex}{}{}
  2,3,5,7,9,11,13,17 et 19 sont les premiers nombres premiers. Ils apparaissent uniquement dans leur propre table de multiplication.
\end{ex}

\begin{rem}{}{}
	Attention, $1$ n'est pas un nombre premier !
\end{rem}


\begin{methode}{Crible d'Erathosthène}{}
La méthode du crible d'Ératosthène permet de dresser la liste des nombres premiers inférieurs ou égaux à $n$.

Nous prendrons l'exemple $n=100$:

\begin{minipage}{.6\textwidth}
  \begin{tabular}{|*{10}{c|}}
\hline
 & 2 & 3 & 4 & 5 & 6 & 7 & 8 & 9 & 10\\
\hline
11 & 12 & 13 & 14 & 15 & 16 & 17 & 18 & 19 & 20\\
\hline
21 & 22 & 23 & 24 & 25 & 26 & 27 & 28 & 29 & 30\\
\hline
31 & 32 & 33 & 34 & 35 & 36 & 37 & 38 & 39 & 40\\
\hline
41 & 42 & 43 & 44 & 45 & 46 & 47 & 48 & 49 & 50\\
\hline
51 & 52 & 53 & 54 & 55 & 56 & 57 & 58 & 59 & 60\\
\hline
61 & 62 & 63 & 64 & 65 & 66 & 67 & 68 & 69 & 70\\
\hline
71 & 72 & 73 & 74 & 75 & 76 & 77 & 78 & 79 & 80\\
\hline
81 & 82 & 83 & 84 & 85 & 86 & 87 & 88 & 89 & 90\\
\hline
91 & 92 & 93 & 94 & 95 & 96 & 97 & 98 & 99 & 100 \\
\hline
\end{tabular}
\end{minipage}
\begin{minipage}{.4\textwidth}
\begin{itemize}
  \item Écrire la liste des entiers de $2$ à $n$.
  \item Éliminer les multiples de $2$, puis de $3$, etc., jusqu'à $\sqrt{n}$.
  \item Les entiers non éliminés sont premiers.
\end{itemize}
\end{minipage}

\end{methode}

\begin{methode}{}{}
  Pour savoir si un nombre $a$ est premier, il suffit de tester s'il est divisible par un nombre premier inférieur à $\sqrt{a}$. En effet, comme les diviseurs peuvent être mis par paire. La \og moitié\fg{}se situe à $\sqrt{a}$.
\end{methode}

\begin{ex}{}{}
  Testons si $\np{7387}$ est premier. Il \og suffit\fg{} de tester les premiers jusqu'à $\sqrt{7387}\approx 85,9$. Donc on teste $2, 3, 5, 7, 11, 13, 17, 19, 23, 29, 31, 37, 41, 43, 47, 53, 59, 61, 67, 71, 73, 79, 83$.

  On trouve qu'il est divisible par 83. Donc il n'est pas premier.
\end{ex}

\begin{histoire}{}{}
	De façon moderne, les nombres premiers sont indispensables en cryptographie pour faire passer des messages codés. Par exemple, ils sont utilisés pour assurer la sécurité d'un paiement bancaire ou la confidentialité d'un appel téléphonique.

	Trouver des grands nombres premiers est un problème difficile et d'actualité. L'Electronic Frontier Foundation offre des prix de calcul coopératif pour encourager les internautes à contribuer à la résolution de problèmes scientifiques par le calcul distribué. Le GIMPS a ainsi reçu 100 000 dollars pour la découverte en 2008 d'un nombre premier avec dix millions de chiffres. C'est même 250 000 dollars qui sont offerts à ceux qui parviendront à trouver un nombre premier avec un milliard de chiffres.
\end{histoire}


\begin{methode}{Déterminer une décomposition en nombres premiers}{}
	On cherche à décomposer le nombre entier $a\in\N$.:
	\begin{enumerate}
		\item On cherche un diviseur $d$ en testant les petits nombres premiers $2,3,5,7,11,13,17,19$.
		\item On écrit alors $a=k\times d$.
		\item Si $k$ est premier alors on a fini. Sinon, on retourne à l'étape 1 en décomposant $k$ cette fois.
	\end{enumerate}
\end{methode}

\begin{ex}{}{}
	Prenons par exemple $84$. On remarque que $84$ est divisible par $2$. On a $84 = 2\times42$. Donc on continue avec $42$.

	On sait que $42$ est divisible par $2$. On a $42=2\times21$. Donc on continue avec $21$.

	On sait que $21=3\times7$. Puisque $3$ et $7$ sont premiers, on a fini.

	On a obtenu la décomposition:
	$84 = 2\times42 = 2\times(2\times21)=2\times2\times(3\times7) = 2^2\times3\times7$
\end{ex}

\begin{theo}{fondamental de l'arithmétique}{}
	Soit $a\in\N$ un entier naturel. Il existe un entier $n\in\N$ des nombres premiers distincts $d_1, d_2, d_3, ..., d_n$ et des exposants $\alpha_1, \alpha_2, \alpha_3, ..., \alpha_n$ tels que:
	$$a=d_1^{\alpha_1}\times d_2^{\alpha_2} \times d_3^{\alpha_3} \times ...\times d_n^{\alpha_n}$$

	On appelle cette écriture une décomposition en nombres premiers.
\end{theo}

\begin{methode}{Utiliser les nombres premiers pour simplifier une fraction $\frac{a}{b}$}{}
   Soient $a\in\Z$ et $b\in\N^*$.

   On décompose $a$ et $b$ en produits de nombres premiers.

   $a = a_1\times a_2\times a_3\times ... \times a_n$ et
   $b = b_1\times b_2\times b_3\times ... \times b_m$

   Après simplification des facteurs, le résultat est une fraction irréductible.
\end{methode}

\begin{ex}{}{}
  $\frac{105}{84} = \frac{3\times 5 \times 7}{2\times 2 \times 3 \times 7} = \frac{5 }{2\times 2} = \frac{5}{4} = \np{1.25}$
\end{ex}

\begin{ex}{Trouver le nombre et l'ensemble des diviseurs de 360}{}
  On décompose $360$ en nombres premiers:
  $$360=2\times180=2\times(2\times90)=2^2\times(2\times45)=2^3\times(3\times15)=2^3\times3\times(3\times5) = 2^3\times3^2\times5$$

Les diviseurs sont alors des sous produits. On peut choisir si on prend le facteur 2 et à quelle puissance entre 1 et 3. On peut choisir si on prend le facteur 3 et à quelle puissance entre 1 et 2. On peut choisir si on prend le facteur 5.

\medskip
Cela donne donc 4 choix pour la puissance de 2. 3 choix pour la puissance de 3 et 2 choix pour la puissance de 5.

Soit un total de $4\times3\times2 = 24$ diviseurs.

L'ensemble des diviseurs est :

$$D = \{2^0\times3^0\times5^0;\quad 2^0\times3^0\times5^1;\quad 2^0\times3^1\times5^0;\quad 2^0\times3^1\times5^1;\quad 2^0\times3^2\times5^0;\quad 2^0\times3^2\times5^1;\quad 2^1\times3^0\times5^0;\quad \dots;\quad 2^3\times3^2\times5^1\}$$
\end{ex}

\newpage
\section*{Exercices}
\sectionexo{Divisibilité}
\begin{bookexos}{47-53 p.66}\end{bookexos}

\sectionexo{Parité}
\begin{bookexo}{73  p.67}\end{bookexo}
\begin{bookexo}{76 p.67}\end{bookexo}
\begin{bookexo}{80 p.67}\end{bookexo}

\sectionexo{Nombres premiers}
\begin{bookexo}{84 p.67}\end{bookexo}
\begin{bookexo}{85 p.67}\end{bookexo}


\sectionexo{Parcours d'approfondissement}
Preuve des règles de divisibilité par 2, 3, 5, 9...


\end{document}